\documentclass[11pt]{article}
\usepackage[margin=1in]{geometry}
\usepackage{booktabs}
\title{VTA, Intact Hemisphere --- Interpretation}
\author{GSE233866, 6-OHDA lesion cohort}
\date{}

\begin{document}
\maketitle

\subsection*{Why this network was built}

Built 2026-08-13, in response to a direct question: the lesion-model analysis had a full
intact/lesioned/combined-DME trio for SNc, but VTA only had a standalone lesioned network
(\texttt{../vta\_lesioned\_network/}), built solely as input to the SNc-vs-VTA-during-lesioning
DME test. There was no VTA-only intact network to pair it with, so no within-VTA
intact-vs-lesioned comparison was possible. This network, and
\texttt{../lesioned\_vs\_intact\_combined\_dme\_vta/}, fill that gap.

\subsection*{Network construction}

4{,}426 cells (the intact/contralateral hemisphere across the same 6 lesion-arm animals used
throughout this section) cluster at soft power 4 ($R^2=0.959$) --- notably lower than every
other lesion-arm network (SNc intact: 6, SNc lesioned: 7, VTA lesioned: 6), and lower than
this project has needed anywhere else except human control tissue. A lower required power
generally reflects a network with less pronounced scale-free structure to begin with.

\subsection*{CACNA1D's position}

CACNA1D falls in the \textbf{pink} module --- at 132 genes, the smallest CACNA1D-containing
module in this project by a wide margin (compare SNc's 855--1{,}640 or even VTA's own other
networks at 316--695) --- with $\mathrm{kME}=0.230$, the lowest standalone-network kME
recorded anywhere in this project.

\subsection*{No significant enrichment --- new information}

Pink returns \textbf{zero} significantly enriched terms across both libraries (826 terms
tested; best result "Morphine addiction," KEGG\_2019\_Mouse, $p_{\mathrm{adj}}=0.120$, well
short of conventional significance). Every other CACNA1D-containing module examined in this
project --- SNc at every stage, and VTA's own healthy-baseline and lesioned networks --- has
at least one term clearing $p_{\mathrm{adj}}<0.02$. This is the first direct evidence in the
project that a CACNA1D module isn't just weakly connected but may not represent a coherent
functional program at all in this particular slice of the data.

\subsection*{Gene-set instability, quantified}

Pink shares only 19 of its 132 genes (14\%) with \texttt{../vta\_lesioned\_network/}'s blue
module (695 genes) --- roughly a third of the equivalent SNc intact-vs-lesioned overlap
(48\%, \texttt{../snc\_intact\_network/} vs.\ \texttt{../snc\_lesioned\_network/}). Combined
with VTA's already-documented pattern (39-module fragmentation at baseline, near-zero
overlap between healthy-baseline and lesioned VTA modules --- see
\texttt{../../healthy\_baseline/vta\_network/INTERPRETATION.tex}), this is now the fourth
independent line of evidence that VTA's networks in this project are considerably less
stable estimates than SNc's, most plausibly because every VTA network here is built from
roughly a quarter to a half of SNc's cell count.

\subsection*{What this means for the region-selectivity question}

This network's low kME, small module size, complete lack of enrichment, and low
cross-network gene overlap should all be read primarily as \textbf{estimate instability},
not as \emph{evidence against} CACNA1D having a role in VTA. A weak, noisy estimate is not
the same as a confirmed null result. The more informative comparison for the
region-selectivity question this network was built to help answer is the paired
kME trajectory: VTA intact (0.230, here) $\to$ VTA lesioned (0.283,
\texttt{../vta\_lesioned\_network/}) --- see
\texttt{../lesioned\_vs\_intact\_combined\_dme\_vta/INTERPRETATION.tex} for that comparison
alongside SNc's, and what it does and does not support.

\end{document}
